\section{The precise interpretation in intuitionistic primitive recursive arithmetic}
\label{sec:arithmetic}
The geometric argument has been arranged so that its exact part is an
arithmetical proof, not an appeal to the totality of an unrestricted real
optimization procedure. This section specifies that statement and the
operations needed to express it.

\subsection{The underlying theory}
We use first-order intuitionistic logic with equality, natural numbers,
zero and successor, a function symbol for each primitive recursive
definition with its defining equations, the successor axioms, and
induction for quantifier-free formulas. This is iPRA in the convention
used by ProveIt. Integers, rationals, finite lists, and matrices below
are natural-number codes; they are not additional logical sorts.
Induction formulas may have parameters. Bounded conjunctions and
disjunctions of decidable numerical conditions are represented by
primitive recursive Boolean functions, so their use does not require
induction for arbitrary quantified formulas.

\subsection{An honest code for the ordered quadratic field}
Encode an integer by a sign and a magnitude, a finite tuple by a fixed
primitive recursive pairing function, and an element of $K=\QQ(\varphi)$
by a triple
\[
 (A,B,D),\qquad A,B\in\mathbb Z,\quad D\in\mathbb N_{>0},
 \qquad \llbracket A,B,D\rrbracket=(A+B\varphi)/D.
\]
Reduced fractions are optional. Equality is checked by cross
multiplication of both coefficients. Addition uses common denominators.
Multiplication before denominator clearing is
\[
 (a+b\varphi)(c+d\varphi)=(ac+bd)+(ad+bc+bd)\varphi.
\]
The conjugate of $a+b\varphi$ is $(a+b)-b\varphi$, so its norm is
$a^2+ab-b^2$. For a nonzero triple, inversion is
\[
 (A,B,D)^{-1}
 =\left(D(A+B),-DB,A^2+AB-B^2\right),
\]
followed by flipping all three signs if the denominator is negative.
These are formulas on integer codes. In particular, we do not assert
that there is an integer satisfying $\varphi^2=\varphi+1$; that would
make an intended encoding vacuous.

For comparison, observe that
\[
 2D\llbracket A,B,D\rrbracket=(2A+B)+B\sqrt5.
\]
Let $c=2A+B$ and $d=B$. If one is zero or both have the same sign,
the sign is immediate. If their signs differ, compare $c^2$ with
$5d^2$: the sign of the field element is
$\operatorname{sign}(c)\operatorname{sign}(c^2-5d^2)$.
This is an exact, decidable test on integers, not a floating-point
test on a radical.

For completeness, $c^2=5d^2$ with integer $c,d$ forces $c=d=0$.
Reduction modulo $5$ first gives $5\mid c$, then $5\mid d$; dividing
by $25$ repeats the equation with smaller absolute values. A finite
descent proof is obtained by bounded search for a least nonzero pair
below a supplied counterexample, or by quantifier-free induction on
the Boolean assertion that no such pair occurs below a given bound.
This proves nonvanishing of the norm and validates the field inverse.
The same elementary sign cases prove compatibility of the order with
addition and multiplication. None requires excluded middle for
uncoded real numbers.

\subsection{The section algorithm and the internal theorem}
Fix the five vertex and edge lists in the order $T,C,O,I,D$. A plane
input consists of four field codes $(n_1,n_2,n_3,t)$, with $n\ne0$.
The algorithm \texttt{SecSq} is defined as follows.
\begin{enumerate}
\item Evaluate the at most twenty heights $h_i=n\cdot v_i-t$.
\item Insert every zero-height original vertex and every intersection
from (\ref{eq:edge}) for a strictly crossed edge. Deduplicate by exact
field equality. The input solids have at most thirty edges.
\item Choose the first nonzero component of $n$ and drop that coordinate.
Sort the projected point list lexicographically by the exact field order.
Use the monotone-chain convex hull scan: delete the last point while the
last turn is nonpositive. This gives the cyclic boundary, with collinear
intermediate points removed. Lift the stored original points back to
three coordinates.
\item Return $\tfrac14\|\sum_iq_i\times q_{i+1}\|^2$ as a field code;
return zero for lists of length at most two.
\end{enumerate}
All invalid codes and the zero-normal case may be totalized to zero;
the theorem uses an explicit validity antecedent. Each loop is bounded
by a fixed number of vertices or edges, and each integer operation has
a standard primitive recursive definition. The convex hull scan's
invariants are orientation signs, orderedness, and containment relative
to finitely many supporting lines. They are Boolean arithmetical
conditions on the code, so their inductive proofs are quantifier free.

The exact upper-bound constants in the unscaled models are
\begin{equation}
 B_T=12,\quad B_C=32,\quad B_O=4,\quad
 B_I=20+32\varphi,\quad B_D=\frac{75}{4}+25\varphi.
                                                               \label{eq:B}
\end{equation}
Let \texttt{leK} be the Boolean comparison of field codes. The
first-order statement proved by the construction is
\begin{equation}
 \forall s,n,t\quad
 \bigl(s<5\ \wedge\ \mathrm{ValidPlane}(n,t)\bigr)
 \ \longrightarrow\
 \mathrm{leK}(\mathrm{SecSq}(s,n,t),B_s)=1.          \label{eq:ipraendpoint}
\end{equation}
All displayed objects here are abbreviations for natural-number
terms and formulas. The body is a decidable arithmetical implication.
The five attaining inputs are the closed codes of the plane equations
in the first table; for each, the stronger equality
$\mathrm{SecSq}(s,n_s,t_s)=B_s$ holds in field-code equality.
Thus attainment is supplied by explicit witnesses, not an existence
principle for a maximum.

\subsection{Why the proof stays within iPRA}
The relevant proof operations can be traced explicitly.
\begin{itemize}
\item \emph{Finite geometry.} Edge clipping, supporting-plane
incidences, orientation tests, affine interpolation, and the shoelace
formula are finite identities and finite lists of inequalities.
Polygon area monotonicity follows from a common finite vertical
subdivision and comparison of endpoint heights. The number of
trapezoids is bounded by the number of polygon vertices.
\item \emph{Centralization.} The reflected-fiber polygon has at most
$4m$ listed vertices for an $m$-vertex slice. Its equal-area statement
is precisely the finite-sum identity (CS5). There is no use of
integration, real completeness, or a general Brunn--Minkowski theorem.
\item \emph{Tetrahedral norm bounds.} Use (\ref{eq:variance}) rather
than introducing a normed-space completion or a square-root function.
All face-vector weights and their slack identities lie in $K$.
\item \emph{Symmetry.} A symmetry is coded by its matrix or its
permutation of the fixed vertex list. The symmetry group is finite,
since three linearly independent vertices determine a linear map.
One may enumerate the images of these three vertices, with at most
$20^3$ candidate triples, retaining the orthogonal maps preserving the
whole vertex list. The finite maximum of a list of field elements is
obtained by an ordinary comparison fold. This supplies the chamber
covering arguments without a maximum theorem for a sphere.
\item \emph{Polynomial positivity.} A polynomial is a finite array
of coefficient/exponent pairs. Addition, convolution, substitution of
linear forms, and evaluation are bounded sums. If every coefficient
is nonnegative and every input variable is nonnegative, evaluation is
nonnegative by induction on a finite prefix sum. Every algebraic
identity used in the certificate is verified by coefficient equality.
\end{itemize}

For the icosahedron, each certificate coefficient is an explicitly
positive integer combination of $1$ and $\sqrt5$. For the dodecahedron,
the single fixed rational enclosure of $\varphi$ in its proof certifies
all coefficient signs. After denominators have been given their
positive signs, every sharp inequality is a finite sum of products
of nonnegative terms. The coefficient tables are therefore concrete
arithmetic proof data, rather than an oracle asserting the output of
a global optimizer.

\paragraph{Strength of the conclusion.}
The preceding text is a human-readable iPRA proof construction: every
induction and every computation has been specified in finitary form.
It does not claim that merely being a terminating Python program
automatically makes a theorem provable in iPRA. The geometric and
arithmetic invariants are the proof of correctness; translating them
into the syntax of \texttt{Prf PRAAx} is the remaining formalization
project described below. In particular, no uniform reflection
principle for iPRA and no universal evaluator for all primitive
recursive programs is assumed.
